% ===== Mục 4.7 =====
\subsubsection{Gom khoản thuốc}
Thí nghiệm S1 đánh giá cách gom khoản trên 176 ảnh tập kiểm tra VAIPE (443 khoản), cắt theo vùng thuốc của nhãn chuẩn.
Tham số được chọn trên 176 ảnh tập kiểm định.
Bảng~\ref{tab:S1} trình bày kết quả.

\begin{table}[H]
\centering
\caption{Gom khoản thuốc trên 176 ảnh tập kiểm tra VAIPE, khoản đúng khi IoU $>0{,}5$ (thí nghiệm S1).}\label{tab:S1}
\small
\begin{tabular}{lrrrr}
\toprule
\textbf{Phương án gom khoản} & \textbf{Precision} & \textbf{Recall} & \textbf{F1} & \makecell[r]{\textbf{IoU}\\\textbf{trung bình}} \\
\midrule
\multicolumn{5}{l}{Tách dòng bằng lược đồ chiếu ngang} \\
\makecell[l]{Phép giãn dọc và luật dòng đậm,\\$s=3$, ngưỡng 1,25 (đề xuất)} & \textbf{0,971} & \textbf{0,975} & \textbf{0,973} & 0,813 \\
(c) Số dòng cố định, $n=2$ & 0,968 & 0,971 & 0,970 & 0,814 \\
(d) Chỉ phép giãn dọc & 0,214 & 0,422 & 0,284 & 0,603 \\
(f) Phép đóng dọc & 0,131 & 0,111 & 0,120 & 0,733 \\
(b) Khe cố định & 0,101 & 0,113 & 0,107 & 0,758 \\
(a) Mỗi dòng một khoản & 0,048 & 0,095 & 0,064 & 0,728 \\
\midrule
\multicolumn{5}{l}{Tách dòng bằng thành phần liên thông} \\
Phương án đề xuất & 0,831 & 0,946 & 0,885 & \textbf{0,833} \\
\bottomrule
\end{tabular}
\end{table}

Trên đơn VAIPE, khe giữa hai dòng trong một khoản xấp xỉ khe giữa hai khoản.
Vì vậy mọi cách chỉ dựa vào khoảng cách, gồm phép giãn, phép đóng và khe cố định, đều đạt F1 không quá $0{,}28$.
Thêm luật ``dòng đậm mở khoản mới'', với độ đậm đo bằng phép co~\eqref{eq:bold}, nâng F1 lên \num{0.973}.
Phương án (c), mỗi khoản 2 dòng, đạt F1 gần bằng (\num{0.970}) chỉ vì mọi đơn VAIPE dùng cùng một mẫu in.
Phương án này không tổng quát cho đơn có khoản 1 dòng hoặc 3 dòng, còn luật dòng đậm không giả định số dòng.
Tách dòng bằng lược đồ chiếu cho F1 cao hơn tách bằng thành phần liên thông (\num{0.973} so với \num{0.885}).
Nguyên nhân là dấu thanh tiếng Việt tạo nhiều thành phần nhỏ lệch tâm dòng.
Khi khoanh vùng thuốc tự động (thí nghiệm S4), IoU vùng đạt \num{0.725} và F1 khoản đầu-cuối đạt \num{0.820}, so với \num{0.219} khi lấy cả trang (Bảng~\ref{tab:S4}, Phụ lục~\ref{app:prescription}).

\subsubsection{Chạy đầu-cuối trên dữ liệu tổng hợp}
\begin{figure}[H]
\centering
\includegraphics[width=0.8\linewidth]{prescription/S2.png}
\caption{F1 khoản thuốc khi chạy đầu-cuối qua các bước (1), (2) và (4) trên tập kiểm tra \bo{5} (thí nghiệm S2).
Mỗi đường chỉ gồm các mẫu bị đúng một loại suy biến.}\label{fig:S2}
\end{figure}

Hình~\ref{fig:S2} cho thấy nghiêng và méo phối cảnh ít ảnh hưởng đến kết quả, với F1 0,85--0,99, và chỉ giảm còn \num{0.69} ở góc 30°.
Nguyên nhân là bước (1) đã nắn ảnh trước khi tách khoản.
F1 chung trên 450 mẫu là \num{0.725}.
Ca lỗi điển hình là bóng đổ, làm F1 giảm từ \num{0.85} xuống \num{0.19}--\num{0.35} dù ảnh nhị phân vẫn đúng (thí nghiệm T3).
Sau khử bóng, nét chữ trong vùng tối mảnh hơn, nên tỉ lệ độ đậm giữa dòng đậm và dòng thường không còn đủ lớn.
Hệ quả là vùng thuốc và ranh giới khoản bị xác định sai.
Mờ mức 3 (F1 \num{0.48}) và dấu đỏ đè lên dòng tên thuốc (F1 \num{0.56}) gây lỗi tương tự.
Hướng sửa là chuẩn hóa độ đậm theo độ tương phản cục bộ của từng dòng.

\subsubsection{Ảnh hưởng lên OCR}
Thí nghiệm S3 đo lợi ích của tách khoản đối với OCR, trên 176 ảnh tập kiểm tra đã soát.
Chữ nhãn chuẩn là các dòng trong vùng thuốc.
Khung bao do bước (4) tự tìm, và Tesseract~5 với gói \texttt{vie} đọc ảnh xám sau bước (2).
Bảng~\ref{tab:S3} trình bày kết quả.

\begin{table}[H]
\centering
\caption{OCR phần thuốc trên 176 ảnh tập kiểm tra VAIPE với các cách cắt vùng khác nhau (thí nghiệm S3).}\label{tab:S3}
\small
\begin{tabular}{lrr}
\toprule
\textbf{Cách cắt vùng trước OCR} & \textbf{CER} & \textbf{F1 từ} \\
\midrule
Cả trang (đối chứng) & 2,815 & 0,365 \\
Vùng thuốc tự động, OCR một lần & \textbf{0,215} & \textbf{0,700} \\
Từng dòng & 0,241 & 0,592 \\
Từng khoản thuốc (pipeline) & 0,225 & 0,658 \\
Từng dòng của nhãn chuẩn (cận trên) & 0,053 & 0,821 \\
\bottomrule
\end{tabular}
\end{table}

OCR cả trang đọc được 87\% số từ nhưng lẫn toàn bộ phần hành chính của đơn, nên CER lên tới \num{2.815}.
Lợi ích chính đến từ khoanh vùng thuốc, giúp CER giảm hơn 12 lần.
OCR theo khoản không cho CER thấp hơn OCR cả vùng thuốc (\num{0.225} so với \num{0.215}), nhưng cho sẵn ranh giới từng thuốc để tách các trường.
Khoảng cách tới cận trên (\num{0.225} so với \num{0.053}) là lỗi của bước khoanh vùng và tách khoản tự động, với F1 \num{0.820} ở thí nghiệm S4.
Trên 5 đơn thật chụp bằng điện thoại (thí nghiệm S5, định tính), chỉ 1 đơn được tách đúng một phần, vì luật dòng đậm gắn với mẫu in VAIPE (Phụ lục~\ref{app:prescription}).
